\begin{table}[!htbp]
\centering
\caption{Sintesis pemenuhan tujuan artefak berdasarkan hasil evaluasi.}
\label{tab:tab6.20_sintesis_pemenuhan_tujuan}
    \begin{tabularx}{\textwidth}{@{}p{0.11\textwidth}L L L@{}}
    \hline
    \textbf{Tujuan} & \textbf{Bukti utama} & \textbf{Temuan} & \textbf{Batas interpretasi} \\
    \hline
    T1 & Prediksi, Clarke Error Grid, finger-stick, conformal & Prediksi layak pada ruang uji; uncertainty mencapai target coverage pada +30/+60 & Belum validasi klinis umum; hipoglikemia tetap menantang \\
    T2 & MRR natural/divergen, cross-fold, oracle & Conditioning unggul kuat pada kasus divergen; tidak universal pada natural & Bergantung pada ketepatan kondisi prediksi dan konfigurasi retrieval \\
    T3 & CDSS workflow, expert evaluation, SUS, fallback & Integrasi doctor-mediated dapat digunakan dan evidence dapat ditelusuri & Usability perlu ditingkatkan; evaluasi ahli n=12; bukan keputusan klinis otomatis \\
    \hline
    \end{tabularx}
\end{table}
